\chapter{Operaciones entre espectros}\label{cap:operaciones}

Hasta ahora cada capa se selecciona y se reconstruye por su cuenta, y los resultados se suman.
Fourierizer ofrece otras tres maneras de combinar las capas, que operan coeficiente por
coeficiente: \emph{Max}, \emph{Cross} y \emph{Product}. Son operaciones no lineales, y es ahí donde
aparecen los sonidos que no se pueden obtener mezclando.

\section{La suma}

En \emph{Sum}, cada capa $i$ se reconstruye en su propia malla y las señales se suman con su
ganancia $g_i$ y su retraso $\tau_i$:
\[
  y[n]=\sum_i g_i\;\mathcal A_i^{\dagger}\bigl(M_i\odot\mathcal A_i x_i\bigr)\bigl[n-\tau_i\Fs\bigr].
\]
Como la inversa es lineal, sumar señales equivale a sumar tablas cuando las capas comparten malla:
$\mathcal A^{\dagger}(Y_1)+\mathcal A^{\dagger}(Y_2)=\mathcal A^{\dagger}(Y_1+Y_2)$. Por eso no hace
falta una malla común, y cada capa puede tener la ventana que mejor le queda. Es la idea original de
la tesis~\cite{rodriguez2025}: \textbf{las selecciones suenan juntas, como capas de una mezcla, pero
recortadas con precisión en tiempo y frecuencia}.

La suma hereda lo bueno y lo malo de mezclar. Si dos selecciones se traslapan sobre contenido
parecido, sus coeficientes se suman como números complejos: se refuerzan si sus fases coinciden y
se cancelan si se oponen. Además, la capa más fuerte tapa a la débil en la zona común, el
\emph{enmascaramiento} del \cref{cap:sonido}.

\section{Una malla común}

Las otras operaciones comparan coeficientes de capas distintas en la \emph{misma} celda $(m,k)$,
así que todas las capas tienen que analizarse en una sola malla. Fourierizer usa Hann con
$N=L=4096$ y $H=512$ (traslape de 88\,\%). Cada capa se coloca primero en la línea de tiempo
común, con su ganancia y su retraso. Se analiza, se multiplica por su máscara (evaluada en
$t-\tau_i$, para que la selección viaje con la capa) y se obtiene su tabla seleccionada
$\chi_i=M_i\odot\mathcal Ax_i$. La operación combina esas tablas en una tabla $Y$, y una sola
inversa da el resultado. Al final, como en la suma, el pico se normaliza a $-0.3$ dBFS.

\section{Max: un collage espectral}

\[
  Y(m,k)=\chi_{i^*}(m,k),\qquad i^*=\arg\max_i\,|\chi_i(m,k)|,
\]
con empates resueltos a favor de la capa de menor número. En cada celda gana la capa más fuerte, con
su magnitud y su fase.

\begin{proposicion}
\emph{Max} es conmutativa, asociativa e idempotente (combinar una capa consigo misma la deja
igual), y $|Y|=\max_i|\chi_i|$. Además, equivale a la suma de las capas con las máscaras binarias
\[
  B_i(m,k)=\begin{cases}1,& i=i^*(m,k),\\ 0,&\text{en otro caso,}\end{cases}
  \qquad Y=\sum_i B_i\odot\chi_i .
\]
\end{proposicion}

\begin{proof}
Las tres propiedades son las del máximo, salvo el desempate, que es fijo y no depende del orden
en que se combinen las capas. La última igualdad es la definición: en cada celda exactamente un
$B_i$ vale $1$.
\end{proof}

La segunda forma explica cómo suena. \emph{Max} es una suma en la que las máscaras las decide el
propio sonido: el plano se reparte entre las capas según quién es más fuerte en cada punto. No hay
enmascaramiento mutuo (cada celda es de una sola capa) ni cancelaciones de fase. El resultado es
más definido que la suma: un \emph{collage} donde cada sonido aparece completo en sus zonas
fuertes. Pero las máscaras $B_i$ son binarias y pueden saltar de celda en celda, y por el
\cref{cap:stft} eso no es consistente. Donde dos capas tienen niveles parecidos y se alternan, la
reconstrucción produce un granulado leve, el \emph{ruido musical} del \cref{cap:mascaras}.

\section{Cross: el cuerpo de uno, el movimiento de otro}

Con dos capas, $A$ y $B$:
\[
  Y=|\chi_A|\;\e^{\iu\angle\chi_B},
\]
y $Y=0$ donde $B$ es nulo (ahí su fase no está definida). Cada celda toma la magnitud de $A$ y la
fase de $B$.

¿Qué guarda la fase? El teorema de desplazamiento (\cref{cap:fourier}) dice que retrasar una señal
$n_0$ muestras multiplica su transformada por $\e^{-\iu\theta n_0}$: \textbf{la posición de los
eventos en el tiempo está escrita en la fase}, como una pendiente en frecuencia. Dentro de cada
fragmento, la fase de $B$ dice \emph{dónde} caen sus golpes y transientes, y cómo evolucionan sus
parciales. La magnitud de $A$ dice \emph{qué} frecuencias suenan y con qué fuerza: su timbre, sus
armónicos, su envolvente espectral.

La tabla $Y$ casi nunca es consistente: la magnitud de un sonido con la fase de otro no es la STFT de
ningún sonido. La reconstrucción la proyecta (\cref{prop:proyeccion}), así que el resultado no tiene
\emph{exactamente} la magnitud de $A$. Tiene la más cercana que se puede sostener con fases parecidas a
las de $B$. Es la misma situación del célebre algoritmo de Griffin y Lim~\cite{griffinlim1984}, que
itera esa proyección para reconstruir un sonido a partir de su magnitud. Fourierizer hace una sola
pasada, y es justo esa tensión entre magnitud y fase la que da el sonido característico de la síntesis
cruzada.

\begin{estudio}
\textbf{Usar Cross.} Pon en $A$ lo que debe aportar el color: un pad, un acorde, una voz sostenida.
Pon en $B$ lo que debe aportar el ritmo y la articulación: un break, una línea de bajo staccato,
lluvia. El resultado es el acorde de $A$ ``tocado'' con la rítmica de $B$. Si $B$ tiene silencios, el
resultado también, porque donde $B$ es nulo no hay fase que prestar.
\end{estudio}

\section{Product: imprimir un sonido sobre otro}

\[
  Y=|\chi_A|\cdot\chi_B .
\]
En cada fragmento, multiplicar el espectro de $B$ por la magnitud de $A$ es \emph{filtrar} $B$ con un
filtro cuya respuesta en frecuencia es $|\chi_A|$. Multiplicar en frecuencia es convolucionar en el
tiempo: el fragmento de $B$ se convoluciona con la respuesta al impulso de fase cero de ese filtro.
El filtro cambia de fragmento en fragmento, siguiendo a $A$. Es la idea del \emph{vocoder}, en el
que un sonido modulador ($A$, típicamente una voz) imprime su envolvente espectral sobre un portador
($B$, típicamente un sintetizador o ruido). La diferencia es que aquí $|\chi_A|$ no se suaviza: $B$
recibe no solo los formantes de $A$, sino \emph{sus armónicos}.

Hay una consecuencia de escala: en decibeles, el nivel de cada celda es la suma de los niveles de
$A$ y de $B$. Los picos se vuelven mucho más picos y los valles más hondos, y el resultado suele ser
más resonante y ``afinado'' que cualquiera de las dos fuentes. La normalización final corrige el
nivel global, no esa forma.

\begin{estudio}
\textbf{Usar Product.} Con $A$ una voz y $B$ un ruido (lluvia, viento, cinta), la lluvia ``canta''
las notas y las vocales de la voz. Con $A$ un acorde y $B$ una batería, la batería queda afinada al
acorde. Como multiplica dos rangos dinámicos, conviene usar selecciones con feather amplio y
revisar el resultado en la vista de análisis: cualquier resonancia aislada se vuelve muy
evidente.
\end{estudio}

\section{Lo que no entró}

Se probaron otras dos operaciones en el demo web, y se descartaron después de escucharlas:
\begin{itemize}
  \item \textbf{Resta espectral}, $\max(|\chi_A|-\beta|\chi_B|,0)\,\e^{\iu\angle\chi_A}$: ``quitarle $B$
        a $A$''. Produce mucho ruido musical, porque la resta deja islas aleatorias de energía en las
        zonas donde las dos magnitudes se parecen.
  \item \textbf{Morph}, una interpolación de magnitudes entre $A$ y $B$ a lo largo del tiempo. Para
        que suene como una transformación y no como un crossfade, hay que interpolar envolventes y
        estructuras armónicas, no coeficientes. Queda como trabajo futuro.
\end{itemize}

\begin{estudio}
\textbf{Qué operación elegir.}
\begin{itemize}
  \item \emph{Sum:} construir un sonido por partes. Es la más predecible y la única que acepta
        mallas distintas por capa y multirresolución.
  \item \emph{Max:} collages entre sonidos completos, con cada uno nítido en sus zonas fuertes.
  \item \emph{Cross:} color de uno con ritmo de otro.
  \item \emph{Product:} imprimir los armónicos de uno sobre la textura de otro.
\end{itemize}
Las tres últimas son no lineales: el resultado depende mucho de los niveles. Mover el nivel de una
capa en el mezclador cambia qué capa gana en \emph{Max} y cuánto pesa cada fuente en \emph{Cross} y
\emph{Product}.
\end{estudio}
